% =============================================================================
\chapter{Content}
\label{ch:content}
% =============================================================================

The Content nav group gathers the cross-project browse surfaces: Avatars,
Library, Media, Scenes, Derived Clips and the Scene Catalogue. Each page
opens to a list / grid with filters and bulk actions. Where detail drill-downs
use a concrete example, we use \textbf{Amouranth} (an avatar inside the
\textbf{SDG} project in the \texttt{x121} pipeline).

% --------------------------------------------------------------------------
\section{Avatars (browse)}
\label{sec:content-avatars}
% Route: /content/avatars
% Component: app/pages/AvatarsPage.tsx

Cross-project browse surface for the platform's primary entity. "Avatar" is the
unified term for what was historically called character / model — one name across
database, backend and UI. Cards click through to the seed-data editor (see the
Page intent callout below).

\screenshot{05-avatars-hero.png}{Cross-project avatars browser.}

\begin{itemize}
  \pagelement{Filter bar}{Project, scene types, tracks, tag filters; free-text search.}
  \pagelement{Avatar grid}{Card per avatar with hero thumbnail, name, project / group, stats.}
  \pagelement{Bulk action bar}{Appears when rows are selected (approve / label / reject).}
  \pagelement{Empty state}{Rendered when filters match zero avatars.}
\end{itemize}

Clicking an avatar card opens the \texttt{SeedDataModal} rather than navigating.
See Chapter~\ref{ch:pipeline-workspace} for the full Avatar detail walk via
the SDG project.

\begin{designnote}
\textbf{Page intent.} The Avatars page is shaped so it could be exposed
externally to partners who supply seed content — the click target is the
seed-data editor, not an internal detail view, and no pipeline-level
controls are surfaced. If we eventually open an external-partner surface,
this is the page we expect to derive it from.
\end{designnote}

% --------------------------------------------------------------------------
\section{Library}
\label{sec:content-library}
% Route: /content/library
% Component: app/pages/LibraryPage.tsx

Studio-level avatar registry that spans projects and pipelines. An approved avatar
prepared for Project A can be imported into Project B without re-running source-image
QA or variant generation.

\screenshot{05-library-hero.png}{Shared avatar library — cross-project view.}

\begin{itemize}
  \pagelement{Filter bar}{Project, group, tag filters; search.}
  \pagelement{Library card grid}{Each card shows avatar stats and project chip.}
  \pagelement{Card}{Clicking opens the \texttt{LibraryAvatarModal} (Amouranth example).}
  \pagelement{Import affordance}{Button to open the Library \texttt{ImportDialog} (UI shown, not submitted).}
\end{itemize}

\screenshotnarrow{05-library-avatar-modal.png}{LibraryAvatarModal — Amouranth profile overview.}

\begin{designnote}
\textbf{Page intent.} The Library is read-only. Its purpose is to browse
the shared pool of avatars across all projects and pipelines without risk
of accidental edits. The only mutating action surfaced here is
\emph{Import into project}, which is an explicit downstream action rather
than an edit of the source data.
\end{designnote}

% --------------------------------------------------------------------------
\section{Media}
\label{sec:content-media}
% Route: /content/media
% Component: app/pages/MediaPage.tsx

Cross-project browser for source images and generated variants — the seed assets
that drive scene generation. The underlying tables (originally image-only) were
renamed to \texttt{source\_media} / \texttt{media\_variants} when video and audio
became first-class generation inputs, which is why the page route is
\texttt{/content/media} rather than \texttt{/content/images}.

\screenshot{05-media-hero.png}{Media browser — default grid view.}
\screenshot{05-media-list.png}{Media browser — list view toggle.}

\begin{itemize}
  \pagelement{View toggle}{Grid / list switch; density adjustable.}
  \pagelement{Filter bar}{Project, avatar, status, track, tag, date.}
  \pagelement{Variant tile / row}{Thumbnail, version, source, approval state, quick-approve / reject.}
  \pagelement{Bulk action bar}{Select all, label, reject, export.}
  \pagelement{Preview}{Clicking a thumbnail opens the \texttt{ImagePreviewModal} (Amouranth example).}
\end{itemize}

\screenshot{05-media-preview-modal.png}{ImagePreviewModal — Amouranth variant.}

% --------------------------------------------------------------------------
\section{Scenes}
\label{sec:content-scenes}
% Route: /content/scenes
% Component: app/pages/ScenesPage.tsx

Cross-project browser over scene video versions — the assembled clips produced by
the recursive generation loop and its re-stitching extensions. Every regeneration
or external import creates a new version; rows here reflect those versions, not
just scenes.

\screenshot{05-scenes-hero.png}{Scenes browser.}

\begin{itemize}
  \pagelement{Filter bar}{Project, avatar, scene type, track, QA status.}
  \pagelement{Scene rows}{Per-scene-version rows with thumb, duration, QA status, actions.}
  \pagelement{Bulk actions}{Approve / reject / label / export.}
  \pagelement{Playback}{Clicking a row opens the \texttt{ClipPlaybackModal} (Amouranth example).}
\end{itemize}

\screenshot{05-scenes-playback-modal.png}{ClipPlaybackModal — Amouranth clip with transport controls.}

% --------------------------------------------------------------------------
\section{Derived Clips}
\label{sec:content-derived-clips}
% Route: /content/derived-clips
% Component: app/pages/DerivedClipsPage.tsx

Derived clips are externally-processed videos — LoRA training chunks, test renders,
reference cuts — that are imported back into the app for review. They carry a
parent-version link, so every derived clip knows which approved clip it was cut
from, and they reuse the same annotation / tag / approval tooling as generated
clips. Bulk import is supported via browser upload and via a server-side directory
scan.

\screenshot{05-derived-clips-hero.png}{Derived clips browser.}

\begin{itemize}
  \pagelement{Grouping}{Derived clips grouped by parent version.}
  \pagelement{Scan \& import}{Scan a directory to register imported clips (UI only).}
  \pagelement{Actions}{Approve / reject / delete / open lineage.}
\end{itemize}

\screenshotnarrow{05-derived-clips-group.png}{A single derived-clip group expanded.}

% --------------------------------------------------------------------------
\section{Catalogue (Scene Catalogue)}
\label{sec:content-catalogue}
% Route: /content/scene-catalogue
% Component: app/pages/SceneCataloguePage.tsx

The Catalogue is the studio-level registry of what can be generated: image types,
scene types, tracks (clothed / topless / speaker / …), per-track workflow and
prompt assignments, and default video settings. It replaces the old flat
\texttt{variant\_applicability} string with a normalized tracks model and unifies
what used to be two separate tables (\texttt{scene\_types} + \texttt{scene\_catalog})
behind a single page.

\begin{designnote}
Three-tier inheritance: catalogue defaults flow into project settings, which can
be overridden per avatar. A scene type can only reference tracks from the same
pipeline, so cross-pipeline mix-ups are rejected at save time.
\end{designnote}

\screenshot{05-catalogue-hero.png}{Scene Catalogue — default tab (Image Types).}

\begin{itemize}
  \pagelement{Seven tabs}{Image Types, Scene Types, Catalogue, Tracks, Workflows, Prompt Defaults, Video Settings.}
\end{itemize}

\subsection{Image Types tab}
Image catalogue — image counterpart to the scene catalogue. Defines generatable
image types (e.g. \texttt{clothed} generated from \texttt{topless} seed) with
workflow assignments and per-track configuration.
\screenshot{05-catalogue-image-types-tab.png}{Image Types tab.}

\subsection{Scene Types tab}
The reusable "recipe" layer: each scene type packages a workflow, LoRA config,
prompt template, and duration — stamped across avatars rather than configured
one-by-one.
\screenshot{05-catalogue-scene-types-tab.png}{Scene Types tab.}

\subsection{Catalogue tab}
Full cross-product view of scene types × tracks — the generation matrix at a
glance. This is the unified view that replaced the split \texttt{scene\_types} /
\texttt{scene\_catalog} pages.
\screenshot{05-catalogue-catalogue-tab.png}{Catalogue tab — full cross-product of scene types and tracks.}

\subsection{Tracks tab}
Tracks are the extensible replacement for the old \texttt{variant\_applicability}
string (clothed / topless / speaker / …). Admins can add, rename, reorder, and
deactivate tracks without schema migrations.
\screenshot{05-catalogue-tracks-tab.png}{Tracks tab.}

\subsection{Workflows tab}
Per-track default workflow assignments. The three-level naming hierarchy
(workflow $\rightarrow$ scene type $\rightarrow$ scene) lets a single workflow
file back multiple scene variants without duplicating entries.

\screenshot{05-catalogue-workflows-tab.png}{Workflows tab — per-track default workflow assignments.}

\subsection{Prompt Defaults tab}
Per-scene-type prompt templates with placeholder substitution from avatar
metadata. Prompt node mapping and additive per-avatar-per-scene overrides let
prompts be edited in-app instead of opening the workflow in ComfyUI.
\screenshot{05-catalogue-prompt-defaults-tab.png}{Prompt Defaults tab.}

\subsection{Video Settings tab}
Default video parameters that feed the recursive generation loop and the
in-workflow chunking strategies (duration, boundary-frame selection,
clothes-off transition behaviour).
\screenshot{05-catalogue-video-settings-tab.png}{Video Settings tab.}

\begin{designnote}
Most rows on these tabs expose inline editors for prompts, workflow
assignments, or settings. We capture tab list views only — editing UI is
shown but never submitted.
\end{designnote}
